\section{An auxiliary cap and its inclusion probabilities}
\label{aux:section}

The determinant obstruction at the main modulus leaves triples with
repeated labels.  We now impose an independent restriction at a second
prime.  This restriction excludes short integer relations among three
distinct residue columns and gives bounds for the probability of each
remaining triple, including triples with equal residue columns.

Retain the main modulus $h\ge 2$, and choose a prime $q$ with
\begin{equation}\label{aux:parameters}
 H=h^2,\qquad h^{100}<q\le 2h^{100},\qquad
 w=\left\lfloor\frac{q}{1000H^2}\right\rfloor.
\end{equation}
Lemma~\ref{norm:prime-interval} supplies such a prime for all sufficiently
large $h$. These choices imply
$q\ge 2000H^2$, and hence
\[
 \frac{q}{2000H^2}\le w\le\frac{q}{1000H^2}.
\]
We identify the integers $0,\ldots,w-1$ with their residues in
$\mathbb F_q$.

The following construction uses an affine part of a classical elliptic
quadric, a cap with no three collinear points; see
Barlotti~\cite[p.~248]{Barlotti1956}. The proof also gives the required
intersection with a small box.

\begin{lemma}\label{aux:cap}
Let $q$ be an odd prime and $1\le w\le q$ an integer.  There is a set
$S\subset\{0,\ldots,w-1\}^3\subset\mathbb F_q^3$ with
$|S|\ge w^3/q$ such that no affine line contains three distinct points
of $S$.
\end{lemma}

\begin{proof}
Choose a nonsquare $\nu\in\mathbb F_q$ and put
\[
 Q(x,y)=x^2-\nu y^2,\qquad
 \Gamma=\{(x,y,Q(x,y)):x,y\in\mathbb F_q\}.
\]
The form $Q$ vanishes only at $(0,0)$.  On a line with nonzero direction
$(v_1,v_2)$ in its first two coordinates, the equation defining
$\Gamma$ is a quadratic equation in the line parameter with nonzero
leading coefficient $Q(v_1,v_2)$.  Such a line meets $\Gamma$ at most
twice.  A line with direction $(0,0,v_3)$ meets $\Gamma$ exactly once.
Thus $\Gamma$, and every translate of it, has the asserted line
property.

For a uniform $b\in\mathbb F_q^3$, every point belongs to
$b+\Gamma$ with probability $|\Gamma|/q^3=1/q$.  Therefore the average
of $|(b+\Gamma)\cap\{0,\ldots,w-1\}^3|$ is $w^3/q$.
Choose a translate attaining at least this average and take its
intersection with the box.
\end{proof}

Fix one such set $S$ for the parameters in \eqref{aux:parameters}, and
write $s=|S|$.  In particular,
\begin{equation}\label{aux:cap-size}
 s\ge\frac{w^3}{q}\ge\frac{q^2}{2000^3H^6}.
\end{equation}
The set $S$ is fixed before any random choices are made.

For $t\in\mathbb F_q$, write $\|t\|_q$ for the least absolute value of
an integer representative of $t$.  Three distinct cap points already
exclude nontrivial relations over $\mathbb F_q$ whose coefficients sum
to zero.  For a shift $a\in\mathbb F_q^3$, a relation among three points
of $a+S$ with integer
coefficients $|x_i|\le H$ and nonzero sum $m=x_1+x_2+x_3$ would force
$\|ma_1\|_q\le3Hw$.  We exclude this possibility by defining the set of
allowed shifts
\begin{equation}\label{aux:allowed-shifts}
 \mathcal A=
 \{a\in\mathbb F_q^3:
       \|ma_1\|_q>3Hw\text{ for every }1\le m\le3H\}.
\end{equation}
Choose $a$ uniformly from $\mathcal A$, and independently choose
$G_q$ uniformly from $\operatorname{GL}_3(\mathbb F_q)$.  Our auxiliary
set is
\begin{equation}\label{aux:set-definition}
 V=G_q(a+S).
\end{equation}
Here and below a triple $u^{(1)},u^{(2)},u^{(3)}$ is
\emph{affinely independent} if
$u^{(2)}-u^{(1)}$ and $u^{(3)}-u^{(1)}$ are linearly independent.
This allows the three columns to have linear rank two.

\begin{lemma}\label{aux:random-set}
Let $H\ge1$ be an integer, let $q\ge2000H^2$ be prime, and put
$w=\lfloor q/(1000H^2)\rfloor$.  Suppose
$S\subset\{0,\ldots,w-1\}^3$ contains no three distinct collinear
points.  The set $\mathcal A$ in \eqref{aux:allowed-shifts} has at least
$q^3/2$ elements.  For every $a\in\mathcal A$ and every invertible
$G_q$, the set $V$ in \eqref{aux:set-definition} has $|S|$ elements,
does not contain zero, and contains no three distinct collinear points.

Moreover, let $x=(x_1,x_2,x_3)\in\mathbb Z^3\setminus\{0\}$ satisfy
$|x|\le H$, where $|x|$ is the Euclidean norm.  For any
$u^{(1)},u^{(2)},u^{(3)}\in V$, the relation
\begin{equation}\label{aux:short-relation}
 x_1u^{(1)}+x_2u^{(2)}+x_3u^{(3)}=0
 \quad\text{in }\mathbb F_q^3
\end{equation}
is impossible if $x_1+x_2+x_3\ne0$.  If $x_1+x_2+x_3=0$, it is
impossible whenever the three columns are affinely independent.
Consequently, \eqref{aux:short-relation} is impossible for three
distinct elements of $V$.
\end{lemma}

\begin{proof}
For each integer $1\le m\le3H<q$, multiplication by $m$ permutes
$\mathbb F_q$.  Since $3Hw<q/2$, a uniform shift violates the
condition for this $m$ with probability $(6Hw+1)/q$.  The union bound
gives
\[
 \Pr(a\notin\mathcal A)
 \le\frac{3H(6Hw+1)}{q}
 \le\frac{18}{1000}+\frac{3}{2000H}<\frac12.
\]
Thus the required uniform choice of $a$ is defined.  Invertible affine
maps preserve cardinality and the line property.  If $a+v=0$ for some
$v\in S$, then $\|a_1\|_q\le w-1\le3Hw$, contrary to the condition
with $m=1$.  Hence $0\notin V$.

Write $u^{(i)}=G_q(a+v^{(i)})$ with $v^{(i)}\in S$, and let
$m=x_1+x_2+x_3$.  If $m\ne0$, then $1\le |m|\le3H$.
Applying $G_q^{-1}$ to \eqref{aux:short-relation} and taking first
coordinates gives
\[
 ma_1=-\sum_{i=1}^3x_i v^{(i)}_1.
\]
The integer on the right, using the representatives
$0\le v^{(i)}_1<w$, has absolute value at most $3H(w-1)$.
Since $\|ma_1\|_q=\||m|a_1\|_q$, this contradicts
\eqref{aux:allowed-shifts}.

If $m=0$, the coefficients in \eqref{aux:short-relation} have sum
zero also in $\mathbb F_q$.  They are not all zero there, since
$x\ne0$ and $|x_i|\le H<q$.  Such a relation contradicts affine
independence: substituting $x_3=-x_1-x_2$ would give a nontrivial
linear relation between $u^{(1)}-u^{(3)}$ and
$u^{(2)}-u^{(3)}$.  Finally, three distinct elements of $V$ are
affinely independent by the line property.
\end{proof}

The restricted shift excludes short relations.  We next show that
this restriction costs only a constant in the inclusion probability
of an affinely independent triple, and obtain the weaker bounds
needed when residue columns coincide.

For an integer $3\times3$ matrix $A$, define
\begin{equation}\label{aux:weight}
 W(A)=\left(\frac{q^3}{s}\right)^3
       \Pr(\text{all columns of }A\bmod q\text{ belong to }V).
\end{equation}
The probability is over $a$ and $G_q$; $S$ remains fixed.

\begin{proposition}\label{aux:weights}
For the distribution in \eqref{aux:set-definition}, with
$s=|S|\ge1$, the following bounds hold with absolute constants:
\begin{equation}\label{aux:weight-bounds}
 \begin{aligned}
 W(A)&\ll1
   &&\text{if the columns of }A\bmod q\text{ are affinely independent},\\
 W(A)&\ll q^3/s
   &&\text{if }\operatorname{rank}(A\bmod q)\ge2,\\
 W(A)&\ll(q^3/s)^2
   &&\text{for every }A.
 \end{aligned}
\end{equation}
If $W(A)>0$, no column of $A\bmod q$ is zero, and its columns are
either affinely independent or include two equal columns.  For an
affinely independent triple with $W(A)>0$, there is no nonzero
$x\in\mathbb Z^3$ with $|x|\le H$ and $Ax=0\bmod q$.
\end{proposition}

\begin{proof}
First replace the allowed shift by an unconditional uniform
$a\in\mathbb F_q^3$.  Under this replacement the map
$v\mapsto G_qv+G_qa$ is uniform in the affine group: conditional on
$G_q$, its translation $G_qa$ is uniform.  Conditioning back on
$a\in\mathcal A$ increases the probability of any event by at most
two, by Lemma~\ref{aux:random-set}.

The affine group acts transitively on the ordered affinely independent
triples in $\mathbb F_q^3$, whose number is
\[
 q^3(q^3-1)(q^3-q).
\]
There are $s(s-1)(s-2)$ such triples in $S$.  Therefore, for an
affinely independent triple of columns,
\[
 W(A)\le
 2\frac{q^9}{s^3}
 \frac{s(s-1)(s-2)}{q^3(q^3-1)(q^3-q)}\ll1.
\]

For the other two bounds, fix any allowed $a$ and put $S_a=a+S$.
If two specified target columns are linearly independent, their
inverse images under a uniform $G_q$ are uniform among the
$(q^3-1)(q^3-q)$ ordered independent pairs.  At most $s^2$ of these
pairs belong to $S_a^2$.  Requiring all three columns to belong to
$G_qS_a$ can only reduce this probability.  Consequently,
\[
 W(A)\le\frac{q^9}{s^3}\frac{s^2}{(q^3-1)(q^3-q)}
 \ll\frac{q^3}{s}
\]
whenever the target matrix has rank at least two.  This estimate
holds for each allowed shift, so also after averaging over $a$.

If any target column is zero, its inclusion probability is zero.
Otherwise fix one target column.  Its inverse image under $G_q$ is
uniform among the $q^3-1$ nonzero vectors, and $S_a$ consists of $s$
nonzero vectors.  Thus
\[
 W(A)\le\frac{q^9}{s^3}\frac{s}{q^3-1}
 \ll\left(\frac{q^3}{s}\right)^2.
\]
This also covers all patterns of equal columns.  The remaining
assertions follow from the line property and short-relation
conclusion in Lemma~\ref{aux:random-set}.
\end{proof}

In the sampling construction, choose $(a,G_q)$ independently of all
labels, all digit choices, and the random change of coordinates at the main modulus.
Given $V$, choose the auxiliary residues of the sampled columns
independently and uniformly from $V$.  Since $|V|=s$ for every
outcome, the probability of any prescribed ordered triple of
auxiliary residues is
\begin{equation}\label{aux:normalized-probability}
 \frac{\Pr(\text{all three residues belong to }V)}{s^3}
 =\frac{W(A)}{q^9},
\end{equation}
where $A$ is any integer lift of that triple.  This identity includes
repeated residues and is unchanged on conditioning on any main-modulus
data.  It is the normalization used when the two residue conditions
are combined by the Chinese remainder theorem.
